\subsection{性质的下降}\label{src:3.1.4-subsection}
\begin{proposition}[拟凝聚层的 fpqc 局部性质]\label{src:3.1.18}
设 $f:S'\to S$ 是概形的 fpqc 态射.
\begin{enumerate}
\item 拟凝聚 $\mathcal O_S$-模的同态 $\mathcal F\to\mathcal G$ 是同构 (或单射, 满射) 当且仅当 $f^*\mathcal F\to f^*\mathcal G$ 是如此.
\item 拟凝聚 $\mathcal O_S$-模 $\mathcal F$ 是有限型的 (或有限表现的, 平坦的, 向量丛, 线丛) 当且仅当 $f^*\mathcal G$ 是如此. 若 $S$ 与 $S'$ 都是 Noether 概形, 则凝聚性也有相同的结论.
\item $S$-概形 $X$ 上的拟凝聚 $\mathcal O_X$-模 $\mathcal F$ 在 $S$ 上平坦, 当且仅当 $\mathcal F$ 到 $X\times_S S'$ 的拉回在 $S'$ 上平坦.
\end{enumerate}
换言之, 上述各性质在 $S$ 上都是 fpqc 局部的.
\end{proposition}
\begin{proof}
(1) 归约到以下代数命题: 若 $A\to A'$ 是忠实平坦的环映射, 则 $A$-模映射 $M\to N$ 是同构 (或单射, 满射) 当且仅当 $M\otimes_A A'\to N\otimes_A A'$ 是如此. 这直接由 $-\otimes_A A'$ 的忠实正合性得到.

(2) 归约到: 若 $A\to A'$ 忠实平坦, 则 $A$-模 $M$ 有限生成 (或有限表现, 平坦, 局部自由且秩为 $r$) 当且仅当 $M\otimes_A A'$ 是如此. 正向蕴涵是显然的. 反之, 若 $M\otimes_A A'$ 有限生成, 取生成元 $x_1',\ldots,x_m'\in M\otimes_A A'$, 写成 $x_i'=\sum_j x_{ij}\otimes a_{ij}'$, 其中 $x_{ij}\in M$, $a_{ij}'\in A'$. 记 $x_{ij}$ 的个数为 $n$, 则映射 $(x_{ij}):A^{\oplus n}\to M$ 是满射, 因为它在沿忠实平坦映射 $A\to A'$ 作基变换后成为满射. 对核重复这一论证, 可知有限表现性质可以下降.

对于平坦性, 假设 $M\otimes_A A'$ 平坦. 由 $A\to A'$ 的忠实平坦性, $M\otimes_A-$ 的正合性等价于 $(M\otimes_A A')\otimes_{A'}(A'\otimes_A-)$ 的正合性, 而后者由 $A\to A'$ 的平坦性及 $A'$-模 $M\otimes_A A'$ 的平坦性得到. 由于局部自由且秩有限等价于有限表现且平坦, 最后的断言也随之成立.

(3) 归约到: 若 $A\to A'$ 忠实平坦, $A\to B$ 是环映射, 则 $B$-模 $N$ 在 $A$ 上平坦当且仅当 $N\otimes_A A'$ 在 $A'$ 上平坦. 这是 (2) 的特例. 另见~\cite[IV2.2.5]{EGA}~与~\cite[Tag 05AY]{SP}.
\end{proof}
下面这个或许有些令人意外的事实, 即正则性沿忠实平坦映射下降, 将会很有用.
\begin{lemma}\label{src:3.1.19}
若 $A\to B$ 是 Noether 局部环之间的平坦局部环映射, 且 $B$ 正则, 则 $A$ 也正则.
\end{lemma}
\begin{proof}
回顾, Noether 局部环 $R$ 正则当且仅当每个有限生成的 $R$-模 $M$ 都有由有限秩自由模组成的有限消解. 而且, 若 $R$ 正则且维数为 $d$, 并且
\begin{equation}\label{src:3.1.20}
0\longrightarrow K\longrightarrow R^{\oplus k_{d-1}}
\longrightarrow\cdots\longrightarrow R^{\oplus k_0}
\longrightarrow M\longrightarrow0
\end{equation}
是 $R$-模的正合序列, 则 $K$ 自由; 见~\cite[定理 $19.12$]{Eis95}~与~\cite[Tag 00OC]{SP}.
若 $M$ 是有限生成的 $A$-模, 选取正合序列~\bookeqref{3.1.20}. 由于 $B$ 正则, $K\otimes_A B$ 自由. 又因局部自由是拟凝聚层的 fpqc 局部性质~(\bookref{3.1.18}(2)), 且 $R$ 是局部环, 故 $K$ 自由. 因此 $A$ 正则. 另见~\cite[IV0.17.3.3]{EGA}~与~\cite[Tag 00OF]{SP}.
\end{proof}
\begin{proposition}[概形性质的 fpqc 下降]\label{src:3.1.21}
设 $X\to Y$ 是概形的 fpqc 态射. 若 $X$ 拟紧 (或局部 Noether, Noether, 整, 约化, 正规, 正则), 则 $Y$ 也如此.
\end{proposition}
\begin{proof}
首先注意, 拟紧性沿任意满射映射下降. 其他各情形归约到以下代数命题: 若 $A\to B$ 是环的忠实平坦映射, 且 $B$ 是 Noether 环 (或整环, 约化环, 正规环, 正则环), 则 $A$ 也如此. 由于 $A\to B$ 忠实平坦, 它是单射, 且对每个理想 $I\subseteq A$ 都有 $I=IB\cap A$. 由 $A\to B$ 的单射性, 整环与约化环的情形是显然的.

对于 Noether 性, 若 $I_1\subseteq I_2\subseteq\cdots$ 是理想升链, 则因 $I_1B\subseteq I_2B\subseteq\cdots$ 稳定, 升链 $I_1=I_1B\cap A\subseteq I_2=I_2B\cap A\subseteq\cdots$ 也稳定.
若 $B$ 是正规整环, 且 $a/b$ 在 $A$ 上整, 其中 $a,b\in A$, 则 $a/b\in B$. 因而 $a\in(b)\subseteq B$, 即 $a:B\to B/bB$ 是零映射. 这一映射是 $a:A\to A/bA$ 的基变换, 故由忠实平坦性知 $a:A\to A/bA$ 是零映射. 因此 $a\in(b)\subseteq A$, 从而 $a/b\in A$.
正则性的断言由引理~\bookref{3.1.19}~得到. 另见~\cite[Tags 033D, 034B, 06QL]{SP}.
\end{proof}
\begin{proposition}[源端性质的 fpqc 下降]\label{src:3.1.22}
设 $X'\to X$ 是概形的 fpqc 态射. 若概形态射 $X\to Y$ 的复合 $X'\to X\to Y$ 光滑 (或平展), 则 $X\to Y$ 光滑 (或平展).
\end{proposition}
\begin{proof}
光滑情形由引理~\bookref{3.1.19}~得到. 平展情形由光滑情形以及如下观察得到: 对 $y\in Y$, 纤维上的映射 $X_y'\to X_y$ 是满射, 因而若 $\dim X_y'=0$, 则 $\dim X_y=0$. 另见~\cite[IV4.17.7.7]{EGA}~与~\cite[Tag 05B5]{SP}.
\end{proof}
但要注意, 光滑性与平展性并不是 fpqc (或 fppf) 局部性质, 即不能说一个性质成立当且仅当它在 fpqc (或 fppf) 覆盖后成立. 例如, 域上存在非光滑 (但必然平坦) 的有限型概形.
\begin{proposition}[概形的 fppf 与光滑局部性质]\label{src:3.1.23}
\begin{enumerate}
\item 若 $X\to Y$ 是概形的 fppf 态射, 则 $X$ 局部 Noether 当且仅当 $Y$ 局部 Noether.
\item 若 $X\to Y$ 是概形的满射光滑态射, 则 $X$ 约化 (或正规, 正则) 当且仅当 $Y$ 也如此.
\end{enumerate}
\end{proposition}
\begin{proof}
正向蕴涵由命题~\bookref{3.1.21}~得到. 对于 (1), 若 $Y$ 局部 Noether, 则由 Hilbert 基定理知 $X$ 也局部 Noether. (2) 归约到以下代数命题: 若 $A\to B$ 是光滑环映射, 且 $A$ 约化 (或正规, 正则), 则 $B$ 也如此. 我们将证明留给读者. 另见~\cite[Tag 034D]{SP}.
\end{proof}
\begin{remark}\label{src:3.1.24}
整环性质不是平展局部的, 例如结点三次曲线有可约的平展邻域 (见例~\bookref{1.5.4}). 约化性, 正规性与正则性都不是 fppf 局部的: 域上存在非约化的有限型概形. 另一方面, 若 $X\to Y$ 是 Noether 概形的平坦态射, $Y$ 正规, 且每个纤维 $X_y$ 都正规, 则 $X$ 正规; 见~\cite[IV2.6.5.4]{EGA}~或~\cite[Tag 0C22]{SP}.
\end{remark}
\begin{remark}\label{src:3.1.25}
若 $A$ 是 Noether 局部环, 则到其完备化的映射 $A\to\widehat A$ 忠实平坦. 若完备化 $\widehat A$ 约化 (或正规, 正则), 则由概形性质的 fpqc 下降~(\bookref{3.1.21}), $A$ 也如此. 逆命题对于正则性成立, 但对于约化性与正规性一般不成立. 不过, 若 $A$ 在域上本质有限型 (或更一般地, 是优良环), 则 $A$ 约化 (或正规) 当且仅当 $\widehat A$ 也如此; 而且此时正规化与完备化可交换. 见~\cite[Tags 07NZ, 0C23]{SP}.
\end{remark}
\begin{proposition}[靶端的 fpqc 局部性质]\label{src:3.1.26}
设 $S'\to S$ 是概形的 fpqc 态射, $P$ 是概形态射的下列性质之一: 满射, 拟紧, 拟分离, 同构, 开浸入, 闭浸入, 单态射, 仿射, 拟仿射, 拟紧局部闭浸入, 局部有限型, 局部有限表现, 分离, 固有, 泛闭, 泛开, 泛浸没, 有限, 局部拟有限, 拟有限, 平坦, fppf, 光滑, 平展, 非分歧或 syntomic. 则 $X\to S$ 具有 $P$ 当且仅当 $X\times_S S'\to S'$ 具有 $P$.
\end{proposition}
换言之, 上述各性质在靶端都是 fpqc 局部的. \cite[第 305--306 页]{Poo17}~中有一张有用的表, 列出概形态射的各种性质及其是否满足 fpqc 下降 (以及是否在基变换, 复合与扩展下保持).
\begin{proof}
为设置记号, 考虑 Cartesian 图
\[
\begin{tikzcd}[column sep=large,row sep=large]
X'\times_X X'\arrow[d]\arrow[dr,phantom,"\square"]\arrow[r,shift left=1.1ex,"p_1"]
\arrow[r,shift right=1.1ex,"p_2"']&
X'\arrow[d,"f'"']\arrow[r,"g'"]\arrow[dr,phantom,"\square"]&
X\arrow[d,"f"]\\
S'\times_S S'\arrow[r,shift left=1.1ex,"p_1"]
\arrow[r,shift right=1.1ex,"p_2"']&S'\arrow[r,"g"']&S
\end{tikzcd}
\]
我们已经知道, 每个性质都在基变换下保持. 不难看出, 满射, 拟紧与拟分离性质可以下降.

对于同构, 设 $h':S'\to X'$ 是 $f'$ 的逆. 由于逆唯一, 作为 $S'\times_S S'\to X'\times_X X'$ 的态射有 $p_1^*h'=p_2^*h'$. 将态射的 fpqc 下降~(\bookref{3.1.10}(2))~应用于 $X'\to X$, $S'$-态射 $h'$ 下降为 $S$-态射 $h:S\to X$. 由此得 $f\circ h=\id_Y$. 为证明 $h\circ f=\id_X$, 注意两个态射 $\id_X,h\circ f:X\to X$ 在预合成 $X'\to X$ 后相等, 故由态射的 fpqc 下降~(\bookref{3.1.7})~知 $h\circ f=\id_X$.

对于开浸入, 若 $f'$ 是开浸入, 则 $f'(X')=g^{-1}(f(X))$ 是开的. 由于 $g$ 泛浸没 (习题~\bookref{A.4.9}), $f(X)$ 是开的, 问题归约到证明 $X\to f(X)$ 是同构, 而这由前一种情形得到. 闭浸入与拟紧局部闭浸入的情形也随之得到.
对于仿射与拟仿射, 可设 $f$ 拟紧且拟分离. 利用与沿 $S'\to S$ 的平坦基变换可交换的典范分解 $f:X\to\Spec_Y f_*\mathcal O_X\to S$, 以及如下事实: $f$ 仿射 (或拟仿射) 当且仅当 $X\to\Spec_Y f_*\mathcal O_X$ 是同构 (或开浸入).

局部有限型 (或局部有限表现) 的情形归约到: 若 $A\to A'$ 忠实平坦, 则环映射 $A\to B$ 有限型 (或有限表现) 当且仅当 $A'\to A'\otimes_B B'$ 是如此.
若 $A'\to A'\otimes_A B$ 有限型, 则有 $A'$-代数生成元 $b_1',\ldots,b_n'$, 可写成 $b_i'=\sum_j a_{ij}'\otimes b_{ij}$, 其中 $a_{ij}'\in A'$, $b_{ij}\in B$. 记由 $b_{ij}$ 生成的 $A$-子代数为 $\widetilde B\subseteq B$. 由于 $\widetilde B\otimes_A A'=B\otimes_A A'$, $A\to A'$ 的忠实平坦性表明 $\widetilde B=B$.
若 $A'\to A'\otimes_A B$ 有限表现, 则由刚才的论证知 $A\to B$ 有限型, 可以写成 $B=A[x_1,\ldots,x_n]/I$. 由于 $A\to A'$ 平坦,
\[
B\otimes_A A'=A'[x_1,\ldots,x_n]/I',
\qquad
I'=I\otimes_{A[x_1,\ldots,x_n]}A'[x_1,\ldots,x_n].
\]
由于 $I'$ 是有限生成理想, 拟凝聚层性质的 fpqc 下降~(\bookref{3.1.18}(2))~表明 $I$ 也有限生成.

平坦的情形很容易, 由命题~\bookref{3.1.18}(2)~得到. 光滑等价于平坦, 局部有限表现且每个纤维光滑, 因而这一情形归约到以下代数事实: 代数闭域 $K$ 上的有限型代数 $A$ 正则当且仅当对每个代数闭域扩张 $L/K$, $A\otimes_K L$ 都正则. 其余情形留给读者. 另见~\cite[IV2.2.6--7 及 IV4.17.7.4]{EGA}~与~\cite[Tag 02YJ]{SP}.
\end{proof}
\begin{proposition}[源端的 fppf, 光滑与平展局部性质]\label{src:3.1.27}
\begin{enumerate}
\item 若 $X'\to X$ 是概形的 fppf 态射, 则概形态射 $X\to Y$ 局部有限表现 (或局部有限型, 满射, 平坦, fppf) 当且仅当 $X'\to X\to Y$ 是如此.
\item 若 $X'\to X$ 是概形的满射光滑态射, 则概形态射 $X\to Y$ 光滑当且仅当 $X'\to X\to Y$ 光滑.
\item 若 $X'\to X$ 是概形的满射平展态射, 则概形态射 $X\to Y$ 平展 (或局部拟有限, 非分歧) 当且仅当 $X'\to X\to Y$ 是如此.
\end{enumerate}
换言之, 各性质在源端是 fppf, 光滑或平展局部的.
\end{proposition}
\begin{proof}
(1) 归约到: 若 $A\to A'$ 是 $R$-代数的忠实平坦有限表现映射, 则 $R\to A$ 有限型 (或有限表现) 当且仅当 $R\to A'$ 是如此. 利用极限方法~(\bookref{B.3.2}~与~\bookref{B.3.3}), 存在有限型 $R$-代数 $A_0$, $R$-代数的忠实平坦有限表现映射 $A_0\to A_0'$, 以及交换图
\[
\begin{tikzcd}[column sep=large,row sep=large]
R\arrow[r]\arrow[dr]&A_0\arrow[r]\arrow[d]&A_0'\arrow[d]\\
&A\arrow[r]&A'\arrow[ul,phantom,"\lrcorner",near start]
\end{tikzcd}
\]
使得 $A'\cong A\otimes_{A_0}A_0'$. 若 $R\to A'$ 有限型, 则必要时扩大 $A_0'$, 可使 $A_0'\to A'$ 满射. 由于 $A\to A'$ 忠实平坦, 这表明 $A_0\to A$ 满射. 因而 $R\to A$ 也有限型. 有限表现情形留给读者.

对于 (2), 由命题~\bookref{3.1.22}, 光滑性事实上沿 fpqc 态射下降. 反之, 光滑态射在复合下保持.
对于 (3), 容易看出局部拟有限性与非分歧性可以下降; 平展性等价于光滑性与非分歧性, 因而平展性也可以下降. 另见~\cite[IV4.17.7.5]{EGA}~与~\cite[Tags 036M, 036T, 036V]{SP}.
\end{proof}
